% SPDX-License-Identifier: MIT
\chapter{Die erste eigene Oberflaeche}\label{chap:first}
\ChapterSummary{An einem kleinen Formular lernen Sie, wie Widgets
erzeugt, besessen und verschachtelt werden. Sie binden einen Button
an eine Aktion und verstehen, welche Aufgaben eine reale
Anwendungshuelle ausserhalb dieses bewusst kleinen Beispiels erledigt.}

\section{Vom Formular zum Komponentenbaum}
Bei SASD UI Toolkit schreiben Sie die sichtbare Oberflaeche als
C++-Objekte. Der \term{Window}-Container ist eine Wurzel,
\term{VBox} ordnet Elemente untereinander an, und \term{Label},
\term{TextField} sowie \term{Button} bilden die sichtbaren Teile.
Eine \term{Container::emplace}-Operation erzeugt eine Komponente,
übernimmt ihr Eigentum und gibt eine Referenz auf sie zurueck.

\begin{lstlisting}[language=SASDCpp,caption={Ein kleiner Formularbaum, bewusst ohne Host/Eventloop}]
#include <sasd/ui/window.hpp>
#include <sasd/ui/vbox.hpp>
#include <sasd/ui/label.hpp>
#include <sasd/ui/text_field.hpp>
#include <sasd/ui/button.hpp>

void buildGreetingForm(sasd::ui::Window& window) {
    using namespace sasd::ui;
    auto& column = window.emplace<VBox>();
    column.setSpacing(1);

    column.emplace<Label>("Wie heissen Sie?");
    auto& name = column.emplace<TextField>();
    auto& message = column.emplace<Label>("Bereit");
    auto& hello = column.emplace<Button>("Begruessen");

    hello.setOnActivated([&name, &message] {
        message.setText("Hallo, " + std::string{name.text()});
    });
}
\end{lstlisting}

Das ist \textbf{ein syntaktisch nutzbarer Formularbaustein},
aber absichtlich \emph{keine vollständige Anwendung}.
Eine laufende Terminalanwendung muss zusaetzlich das
\term{TerminalBackend} öffnen, Eingaben abholen,
Fokus und Pointer-Routing betreiben, Layout ausführen
und die Darstellung aktualisieren. Diese Verantwortung
zeigt das ausführlich kommentierte
\filepath{examples/terminal_form_demo.cpp}; es ist der
richtige Ausgangspunkt für einen eigenen Host.

\section{Was die paar Zeilen wirklich tun}
Die Referenz \texttt{column} ist kein Besitzer.
Die \texttt{Window}-Wurzel besitzt ihre \texttt{VBox},
und die \texttt{VBox} besitzt wiederum ihre Kinder.
Die Referenzen \texttt{name}, \texttt{message} und
\texttt{hello} sind gültig, solange die Komponenten im
Baum verbleiben. Beim Freigeben des Fensters endet auch
ihre Lebensdauer.

Der Lambda-Callback liest den aktuellen Text des
\term{TextField} und ersetzt die Beschriftung des Labels.
Die sichtbare Aktualisierung wird durch die
Invalidierungs-/Praesentationsmechanik der laufenden
Anwendung verarbeitet. Der Callback allein öffnet
keine Terminal-Session und pumpt keine Ereignisse.

\section{Eine naheliegende Erweiterung}
Ein zweiter Button kann die Eingabe zuruecksetzen:
\begin{lstlisting}[language=SASDCpp]
auto& clear = column.emplace<Button>("Leeren");
clear.setOnActivated([&name, &message] {
    name.setText("");
    message.setText("Bereit");
});
\end{lstlisting}
Die gezeigten Zeilen gehören \emph{innerhalb} von
\texttt{buildGreetingForm}, nach dem Aufbau der zuvor
verwendeten Variablen. Die Reihenfolge ist bedeutsam:
Sie müssen den referenzierten Widgets längeres Leben
garantieren als den installierten Callbacks.

\section{Der nächste Schritt: ein vollständiger Host}
In der realen Terminal-Demo folgen auf den Widget-Aufbau
die Einrichtung von \term{FocusManager}, \term{PointerRouter},
Praesentations-Sink und den Callback-Bindungen.
Der Host pollt Events, vermittelt sie an die
zuständigen Controls und praesentiert den geaenderten
Frame. Das ist kein Versehen und keine fehlende
magische \texttt{Application::run}-Methode:
\term{Application} bietet derzeit unter anderem
\texttt{processEvents()} und
\texttt{processRoutedEvents()}, aber keinen versteckten
universellen Desktop-Message-Loop.

\section{Was Sie aus dem Beispiel mitnehmen}
Der Anwendungsautor beschreibt \emph{Bedeutung und Verhalten}.
Er verwaltet eine eindeutige Eigentumshierarchie und
reagiert auf Ereignisse. Backend, Eingabegerate, Messung
und Darstellung werden bewusst als separate Teile
komponiert, damit die gleiche Semantik später auf
mehreren Ausgabegeraten tragfaehig bleibt.

\PracticalTip{Ersetzen Sie zuerst in der Terminal-Demo nur
den Formulartitel und einen Button-Callback. Vergleichen Sie
danach, was beim Hinzufuegen eines Widgets an Fokusfolge,
Layout und Darstellung ebenfalls betroffen ist.}
